\section{Point-to-Point Protocol}
    \begin{frame}[fragile]{Point-to-Point Protocol (PPP)}
        \begin{itemize}
            \item powstał w 1994
            \item auto-negocjacja prędkości łacza, wielkosci buforów itp
            \item prosta serializacja/deserializacja
            \item wykrywanie błędów i retransmisja
            \item mechanizmy kompresji i szyfrowania
            \item PPP over Serial (PPPoS)
            \item PPP over Ethernet (PPPoE)
            \item PPP over ATM (PPPoA, PPPoEoA)
            \item współcześnie używany przez modemy GSM (AT+PPP) \\ oraz dostawców internetu typu Orange
        \end{itemize}
    \end{frame}

    \begin{frame}[fragile]{Point-to-Point Protocol (PPP)}{c.d.}
        \begin{itemize}
            \item urządzenie "server" czeka na rozpoczęcie sesji
            \item urządzenie "client" rozpoczyna sesję
            \item negocjacja parametrów łącza
            \item urządzenie "server" ustala IP po obu stronach
            \item podsieć gotowa do użytku
        \end{itemize}

        \begin{lstlisting}[language=bash]
sudo pppd /dev/ttyUSB0 921600 192.168.7.1:192.168.7.2 local noauth debug nodetach nocrtscts
        \end{lstlisting}
    \end{frame}
